\chapter{Locally Covariant QFT and Hadamard States}
\label{ch:locally-covariant-qft-hadamard}

\ConventionsMixed

Curved-spacetime QFT is a background-field problem.  The Lorentzian metric,
orientation, time orientation, spin structure, background gauge bundle, and
external sources are held fixed, while the quantum observables are assigned
locally and covariantly to each background.  The absence of a preferred
translation-invariant vacuum is replaced by a condition on states: their
short-distance singularity must have the same positive-frequency geometric
form as the Minkowski vacuum when examined in a sufficiently small convex
normal neighborhood.

This chapter fixes the functorial language and then constructs the free
Klein--Gordon example in detail.  The example is important because the
time-slice axiom, causal propagation, commutation relations, and Hadamard
subtraction are visible in one model without using a perturbative expansion.

\section{Spacetimes And Causal Embeddings}

Let \(M\) be a smooth \(D\)-manifold with Lorentzian metric \(g\) of
mostly-plus signature.  A tangent vector \(v\in T_xM\) is causal when
\(g_x(v,v)\le0\), timelike when \(g_x(v,v)<0\), and null when
\(g_x(v,v)=0\) and \(v\ne0\).  A time orientation is a continuous choice of
future cone in the causal cone at every point.  For a subset \(S\subset M\),
write \(J_M^\pm(S)\) for the causal future and causal past of \(S\).

\begin{definition}[Globally hyperbolic spacetime]
\label{def:globally-hyperbolic-spacetime}
An oriented and time-oriented Lorentzian manifold \((M,g)\) is globally
hyperbolic when it has no closed causal curve and when
\[
  J_M^+(p)\cap J_M^-(q)
\]
is compact for every \(p,q\in M\).  A Cauchy surface is a subset
\(\Sigma\subset M\) intersected exactly once by every inextendible timelike
curve.  For globally hyperbolic \(M\), smooth spacelike Cauchy surfaces
exist, and \(M\) is diffeomorphic to \(\mathbb R\times\Sigma\).
\end{definition}

The compactness condition is the analytic input behind finite propagation
statements.  If \(K\subset M\) is compact, then
\[
  J_M^+(K)\cap J_M^-(K')
\]
is compact for every compact \(K'\subset M\).  This property is used below
when a solution is cut off between two Cauchy surfaces.

\begin{definition}[Causally convex subset]
\label{def:causally-convex-subset}
An open subset \(O\subset M\) is causally convex when every causal curve in
\(M\) whose endpoints lie in \(O\) is entirely contained in \(O\).
\end{definition}

Causal convexity is the condition that lets a field equation be restricted
to a subspacetime without importing causal propagation from the complement.

\begin{definition}[Background category]
\label{def:background-category-loc}
Let \({\bf Loc}\) be the category defined as follows.

An object is an oriented, time-oriented, globally hyperbolic Lorentzian
manifold \((M,g)\), possibly equipped with fixed background bundles and
smooth background sources.  A morphism
\[
  \chi:(M,g)\longrightarrow (N,h)
\]
is a smooth embedding with open causally convex image such that
\[
  \chi^*h=g
\]
and such that \(\chi\) preserves the declared orientation, time orientation,
and background bundle/source data.  A morphism \(\chi:M\to N\) is a Cauchy
morphism when \(\chi(M)\) contains a Cauchy surface of \(N\).
\end{definition}

Composition of morphisms is ordinary composition of embeddings.  The image
of a composite is causally convex: if a causal curve in the final spacetime
has endpoints in the composite image, causal convexity of the second image
first keeps the curve in the intermediate image, and causal convexity of the
first image then keeps its pullback in the original image.  Thus the above
data form a category.

\section{Algebra-Valued Functors}

Let \({\bf Alg}_*\) denote a category of unital complex \(\ast\)-algebras.
Its morphisms are unit-preserving injective \(\ast\)-homomorphisms.  In
analytic versions one uses locally convex topological \(\ast\)-algebras and
continuous morphisms; the algebraic definition is enough for the functorial
axioms in this chapter.

\begin{definition}[Locally covariant QFT]
\label{def:locally-covariant-qft}
A locally covariant QFT is a covariant functor
\[
  \mathfrak A:{\bf Loc}\longrightarrow{\bf Alg}_*
\]
with the following additional properties.

\begin{enumerate}[label=(\roman*)]
\item If \(\chi:M\to N\) is a Cauchy morphism, then
\[
  \mathfrak A(\chi):\mathfrak A(M)\longrightarrow\mathfrak A(N)
\]
is an isomorphism.  This is the time-slice axiom.

\item If \(\chi_i:M_i\to N\), \(i=1,2\), have causally disjoint images, then
\[
  [\,\mathfrak A(\chi_1)(A_1),\mathfrak A(\chi_2)(A_2)\,]=0
\]
for all \(A_i\in\mathfrak A(M_i)\).  This is Einstein causality.
\end{enumerate}
\end{definition}

The injectivity of morphisms in \({\bf Alg}_*\) records isotony: an
admissible spacetime inclusion identifies the algebra of the smaller
background with a subalgebra of the larger one.

\begin{definition}[Kinematic local algebra]
\label{def:kinematic-local-algebra}
Let \(O\subset M\) be open, causally convex, and globally hyperbolic with
the induced metric and orientations.  Let
\(\iota_{M;O}:O\hookrightarrow M\) be the inclusion morphism in
\({\bf Loc}\).  The kinematic algebra of \(O\) inside \(M\) is
\[
  \mathfrak A^{\rm kin}(M;O)
  =
  \mathfrak A(\iota_{M;O})(\mathfrak A(O))
  \subset \mathfrak A(M).
\]
\end{definition}

\begin{remark}[Local net associated to a locally covariant theory]
\label{rem:kinematic-net-from-lcqft}
For each fixed \(M\), the assignment
\[
  O\longmapsto \mathfrak A^{\rm kin}(M;O)
\]
on causally convex globally hyperbolic open subsets is isotone and satisfies
Einstein causality.
\end{remark}

If \(O_1\subset O_2\), the inclusion factors as
\[
  O_1\stackrel{\iota_{O_2;O_1}}{\longrightarrow}O_2
  \stackrel{\iota_{M;O_2}}{\longrightarrow}M .
\]
Functoriality gives
\[
  \mathfrak A(\iota_{M;O_1})
  =
  \mathfrak A(\iota_{M;O_2})\circ
  \mathfrak A(\iota_{O_2;O_1}),
\]
so the image for \(O_1\) is contained in the image for \(O_2\).  If
\(O_1\) and \(O_2\) are causally disjoint in \(M\), the two inclusion
morphisms into \(M\) satisfy the hypothesis of Einstein causality in
Definition~\ref{def:locally-covariant-qft}, and the commutator of the two
image algebras vanishes.

Locally covariant fields are defined by naturality.  Let
\({\bf Test}_E\) be the functor assigning to \(M\) the vector space
\(\Gamma_c(M,E_M)\) of compactly supported smooth sections of a natural
bundle \(E_M\), and to a morphism \(\chi:M\to N\) the pushforward
\(\chi_*:\Gamma_c(M,E_M)\to\Gamma_c(N,E_N)\).

\begin{definition}[Locally covariant field]
\label{def:locally-covariant-field}
A locally covariant field of type \(E\) is a natural transformation
\[
  \Phi:{\bf Test}_E\Longrightarrow \mathfrak A .
\]
Equivalently, for every \(M\) and every \(f\in\Gamma_c(M,E_M)\) there is an
element \(\Phi_M(f)\in\mathfrak A(M)\), and every morphism
\(\chi:M\to N\) satisfies
\[
  \mathfrak A(\chi)\Phi_M(f)=\Phi_N(\chi_*f).
  \label{eq:locally-covariant-field-naturality}
\]
\end{definition}

The condition is a commutative square:
\[
\begin{tikzcd}[column sep=large,row sep=large]
  \Gamma_c(M,E_M) \arrow[r,"\Phi_M"] \arrow[d,"\chi_*"']
  & \mathfrak A(M) \arrow[d,"\mathfrak A(\chi)"]\\
  \Gamma_c(N,E_N) \arrow[r,"\Phi_N"']
  & \mathfrak A(N).
\end{tikzcd}
\]

\section{States As Contravariant Data}

A state on a unital \(\ast\)-algebra \(\mathcal A\) is a linear functional
\(\omega:\mathcal A\to\mathbb C\) satisfying
\[
  \omega(\mathbf 1)=1,
  \qquad
  \omega(A^*A)\ge0
  \quad\hbox{for all }A\in\mathcal A .
\]
Every morphism \(\alpha:\mathcal A\to\mathcal B\) pulls states on
\(\mathcal B\) back to states on \(\mathcal A\) by composition:
\[
  \alpha^*\omega=\omega\circ\alpha .
\]
Thus a state assignment for a locally covariant theory is contravariant.

\begin{definition}[Locally covariant state space]
\label{def:locally-covariant-state-space}
A locally covariant state space for \(\mathfrak A\) is a choice of subsets
\[
  \mathcal S(M)\subset \{\hbox{states on }\mathfrak A(M)\}
\]
such that for every morphism \(\chi:M\to N\),
\[
  \omega\in\mathcal S(N)
  \quad\Longrightarrow\quad
  \omega\circ\mathfrak A(\chi)\in\mathcal S(M).
\]
\end{definition}

Hadamard states provide the state spaces used for free scalar fields on
curved backgrounds.  Their definition is local in spacetime and stable under
restriction to causally convex subregions.

\section{The Free Klein--Gordon Functor}

Fix real constants \(m^2\ge0\) and \(\xi\in\mathbb R\).  For every
\((M,g)\in{\bf Loc}\), define the real differential operator
\[
  P_M=-\nabla^\mu\nabla_\mu+m^2+\xi R[g]
\]
on \(C^\infty(M,\mathbb R)\).  It is formally self-adjoint with respect to
the pairing
\[
  (f,h)_M=\int_M f h\,\dd\operatorname{vol}_g .
\]
Global hyperbolicity gives unique retarded and advanced Green operators
\[
  E_M^{\rm ret},E_M^{\rm adv}:C_c^\infty(M)\longrightarrow C^\infty(M)
\]
with
\[
  P_M E_M^{\rm ret}f=f,\qquad
  P_M E_M^{\rm adv}f=f,
\]
\[
  E_M^{\rm ret}P_M f=f,\qquad
  E_M^{\rm adv}P_M f=f,
\]
and support restrictions
\[
  \operatorname{supp}E_M^{\rm ret}f\subset J_M^+(\operatorname{supp}f),
  \qquad
  \operatorname{supp}E_M^{\rm adv}f\subset J_M^-(\operatorname{supp}f).
\]
The causal propagator is
\[
  E_M=E_M^{\rm ret}-E_M^{\rm adv}.
\]

\begin{lemma}[Kernel of the causal propagator]
\label{lem:kg-causal-propagator-kernel}
For \(f\in C_c^\infty(M)\),
\[
  E_M f=0
  \quad\Longleftrightarrow\quad
  f=P_M h\hbox{ for some }h\in C_c^\infty(M).
\]
\end{lemma}

\begin{proof}
If \(f=P_Mh\) with compactly supported \(h\), then
\[
  E_M f=E_M P_M h=h-h=0,
\]
using the defining identities for retarded and advanced Green operators.

Conversely assume \(E_Mf=0\).  Then
\[
  u:=E_M^{\rm ret}f=E_M^{\rm adv}f .
\]
The retarded expression has support in \(J_M^+(\operatorname{supp}f)\), and
the advanced expression has support in \(J_M^-(\operatorname{supp}f)\).
Thus
\[
  \operatorname{supp}u
  \subset
  J_M^+(\operatorname{supp}f)\cap J_M^-(\operatorname{supp}f),
\]
which is compact by global hyperbolicity.  Finally \(P_Mu=f\).  Hence
\(f\in P_M C_c^\infty(M)\).
\end{proof}

Define the real vector space of smeared solutions by
\[
  \mathcal V(M)
  =
  C_c^\infty(M,\mathbb R)/P_M C_c^\infty(M,\mathbb R),
  \qquad
  [f]_M \hbox{ for the class of }f .
\]
Define
\[
  \sigma_M([f]_M,[h]_M)
  =
  \int_M f\,E_M h\,\dd\operatorname{vol}_g .
  \label{eq:kg-symplectic-form}
\]
This is well-defined because \(E_MP_M=0\) on compactly supported test
functions and because \(P_M\) is formally self-adjoint.  It is antisymmetric:
\[
  \int_M f\,E_Mh\,\dd\operatorname{vol}_g
  =
  -\int_M h\,E_Mf\,\dd\operatorname{vol}_g ,
\]
which follows by applying formal self-adjointness to retarded and advanced
Green operators and using
\[
  (E_M^{\rm ret})^*=E_M^{\rm adv}.
\]
Lemma~\ref{lem:kg-causal-propagator-kernel} implies that
\(\sigma_M\) is nondegenerate on \(\mathcal V(M)\).

\begin{definition}[CCR algebra of the scalar field]
\label{def:curved-kg-ccr-algebra}
The algebra \(\mathfrak A_{\rm KG}(M)\) is the unital \(\ast\)-algebra
generated by symbols \(\Phi_M([f]_M)\), linear in \([f]_M\), with relations
\[
  \Phi_M([f]_M)^*=\Phi_M([f]_M),
  \qquad
  [\Phi_M([f]_M),\Phi_M([h]_M)]
  =
  \ii\,\sigma_M([f]_M,[h]_M)\mathbf 1 .
\]
\end{definition}

The quotient by \(P_MC_c^\infty(M)\) is the equation of motion:
\[
  \Phi_M(P_Mh)=0 .
\]
The commutator is determined by the causal propagator and therefore by the
hyperbolic equation, not by a choice of state.

\begin{proposition}[Functoriality of the free scalar field]
\label{prop:kg-functoriality}
For every morphism \(\chi:M\to N\), the pushforward of test functions
induces a symplectic linear map
\[
  \mathcal V(\chi):\mathcal V(M)\longrightarrow\mathcal V(N),
  \qquad
  [f]_M\longmapsto [\chi_*f]_N .
\]
Consequently there is an injective \(\ast\)-homomorphism
\[
  \mathfrak A_{\rm KG}(\chi):\mathfrak A_{\rm KG}(M)\to
  \mathfrak A_{\rm KG}(N)
\]
defined by
\[
  \mathfrak A_{\rm KG}(\chi)\Phi_M([f]_M)
  =
  \Phi_N([\chi_*f]_N).
\]
\end{proposition}

\begin{proof}
The embedding is isometric and preserves background scalar data, so
\[
  P_N\chi_*f=\chi_*P_Mf .
\]
Therefore \(\chi_*\) maps \(P_MC_c^\infty(M)\) into
\(P_NC_c^\infty(N)\), giving a map on quotients.  Causal convexity gives
compatibility of causal propagators:
\[
  E_N\chi_*f=\chi_*E_Mf .
  \label{eq:causal-propagator-embedding-compatibility}
\]
Indeed, both sides solve \(P_Nu=\chi_*f\) in \(\chi(M)\), both vanish outside
the same retarded-minus-advanced causal support determined by
\(\chi(\operatorname{supp}f)\), and uniqueness of retarded and advanced
solutions gives equality.  The volume form is preserved by \(\chi\), hence
\[
  \sigma_N([\chi_*f]_N,[\chi_*h]_N)
  =
  \int_N \chi_*f\,E_N\chi_*h\,\dd\operatorname{vol}_h
  =
  \int_M f\,E_Mh\,\dd\operatorname{vol}_g
  =
  \sigma_M([f]_M,[h]_M).
\]
The universal property of the CCR algebra then gives a
\(\ast\)-homomorphism preserving the defining commutator.  Injectivity
follows from injectivity of the induced symplectic map; if
\([\chi_*f]_N=0\), then \(E_N\chi_*f=0\), so
\(\chi_*E_Mf=0\), hence \(E_Mf=0\) and \([f]_M=0\) by
Lemma~\ref{lem:kg-causal-propagator-kernel}.
\end{proof}

\begin{proposition}[Time-slice property for the free scalar]
\label{prop:kg-time-slice}
If \(\chi:M\to N\) is a Cauchy morphism, then
\[
  \mathfrak A_{\rm KG}(\chi):\mathfrak A_{\rm KG}(M)\to
  \mathfrak A_{\rm KG}(N)
\]
is an isomorphism.
\end{proposition}

\begin{proof}
It suffices to prove that
\(\mathcal V(\chi):\mathcal V(M)\to\mathcal V(N)\) is surjective, since
Proposition~\ref{prop:kg-functoriality} already gives injectivity and the
CCR algebra is generated by the smeared fields.

Let \(f\in C_c^\infty(N)\) and set \(u=E_Nf\).  Choose two smooth spacelike
Cauchy surfaces \(\Sigma_-\) and \(\Sigma_+\) contained in \(\chi(M)\) such
that \(\Sigma_+\) lies to the future of \(\Sigma_-\) and the compact set
\[
  J_N(\operatorname{supp}f)\cap J_N^+(\Sigma_-)\cap J_N^-(\Sigma_+)
\]
contains the part of \(u\) between the two surfaces.  Choose a smooth
function \(\eta:N\to[0,1]\) that is \(0\) to the past of \(\Sigma_-\), is
\(1\) to the future of \(\Sigma_+\), and whose derivative is supported in
\(\chi(M)\).  Define
\[
  g=P_N(\eta u).
\]
Since \(P_Nu=0\), the support of \(g\) lies where derivatives of \(\eta\)
are nonzero, hence \(g\in C_c^\infty(\chi(M))\).  Set
\[
  w=\eta u .
\]
The section \(w\) vanishes to the past of \(\Sigma_-\) and satisfies
\(P_Nw=g\), so uniqueness of the retarded Green solution gives
\[
  E_N^{\rm ret}g=w .
\]
The section \(w-u=(\eta-1)u\) vanishes to the future of \(\Sigma_+\) and
also satisfies \(P_N(w-u)=g\), so uniqueness of the advanced Green solution
gives
\[
  E_N^{\rm adv}g=w-u .
\]
Therefore
\[
  E_Ng=E_N^{\rm ret}g-E_N^{\rm adv}g=w-(w-u)=u=E_Nf .
\]
Thus \(E_N(g-f)=0\).
By Lemma~\ref{lem:kg-causal-propagator-kernel}, \(g-f=P_Nk\) for some
\(k\in C_c^\infty(N)\).  Thus \([f]_N=[g]_N\).  Since
\(\operatorname{supp}g\subset\chi(M)\), there is
\(\widetilde g\in C_c^\infty(M)\) with \(g=\chi_*\widetilde g\), and
\([f]_N=\mathcal V(\chi)[\widetilde g]_M\).
\end{proof}

\section{Quasifree States}

A state \(\omega\) on \(\mathfrak A_{\rm KG}(M)\) has \(n\)-point
distributions
\[
  \omega_n(f_1,\ldots,f_n)
  =
  \omega\!\left(\Phi_M([f_1])\cdots\Phi_M([f_n])\right).
\]
It is quasifree with vanishing one-point function when all odd \(n\)-point
distributions vanish and all even \(n\)-point distributions are sums over
pairings:
\[
  \omega_{2n}(f_1,\ldots,f_{2n})
  =
  \sum_{\hbox{\scriptsize pairings }P}
  \prod_{\{i,j\}\in P}\omega_2(f_i,f_j).
\]
The two-point distribution \(\omega_2\in\mathcal D'(M\times M)\) must
satisfy
\[
  P_{M,x}\omega_2=0,\qquad P_{M,y}\omega_2=0,
\]
\[
  \omega_2(f,h)-\omega_2(h,f)=\ii\,\sigma_M([f],[h]),
\]
and positivity:
\[
  \omega_2(\bar f,f)\ge0
  \quad\hbox{for all }f\in C_c^\infty(M,\mathbb C).
\]
Conversely, a bidistribution with these properties defines a quasifree state
on the CCR algebra.

\section{Hadamard Form}

Let \(x\) and \(y\) lie in a convex normal neighborhood \(U\subset M\).  Let
\[
  \sigma(x,y)
\]
be Synge's world function: one half of the signed squared geodesic distance
from \(x\) to \(y\).  Let \(T\) be any smooth time function on \(U\), and set
\[
  \sigma_\epsilon(x,y)
  =
  \sigma(x,y)+\ii\epsilon\bigl(T(x)-T(y)\bigr)+\epsilon^2,
  \qquad \epsilon>0 .
\]
The limit \(\epsilon\downarrow0\) is taken in the distributional sense.

In four spacetime dimensions, the local Hadamard singular distribution for
\(P_M\) has the form
\begin{equation}
  H_{\epsilon,\mu}(x,y)
  =
  \frac{1}{8\pi^2}
  \left[
    \frac{U(x,y)}{\sigma_\epsilon(x,y)}
    +V(x,y)\log\!\bigl(\mu^2\sigma_\epsilon(x,y)\bigr)
  \right],
  \label{eq:hadamard-parametrix-local}
\end{equation}
where \(\mu>0\) is a length inverse.  The coefficient \(U\) is
\[
  U(x,y)=\Delta(x,y)^{1/2},
\]
with \(\Delta\) the van Vleck determinant.  It is characterized by the
transport equation
\[
  2\sigma^{;\mu}\nabla_\mu U
  +
  (\Box_g\sigma-D)U=0,
  \qquad
  U(x,x)=1 .
\]
The function \(V\) is locally determined by \(g,m,\xi\) through the
requirement that \(P_{M,x}H_{\epsilon,\mu}(x,y)\) have no singular part away
from the diagonal.  Equivalently, expanding
\[
  V(x,y)\sim\sum_{n\ge0}V_n(x,y)\sigma(x,y)^n
\]
gives transport equations along the geodesic from \(y\) to \(x\), with
initial data fixed recursively by \(P_M\) and the previous coefficients.
The smooth part of a state is not included in \(H_{\epsilon,\mu}\).

\begin{definition}[Hadamard quasifree state]
\label{def:hadamard-quasifree-state}
A quasifree state \(\omega\) of the Klein--Gordon field is Hadamard when, in
every convex normal neighborhood,
\[
  \omega_2(x,y)-H_{\epsilon,\mu}(x,y)
\]
has a smooth distributional limit as \(\epsilon\downarrow0\).  The resulting
smooth function is denoted \(W_\omega(x,y)\).
\end{definition}

Changing the time function \(T\) in \(\sigma_\epsilon\) changes only the
boundary-value prescription used to define the same positive-frequency
singularity.  Changing \(\mu\) changes the subtraction by
\[
  \frac{1}{8\pi^2}V(x,y)\log(\mu'^2/\mu^2),
\]
which is smooth near the diagonal because \(V\) is smooth.

\begin{remark}[State-independent singular part]
\label{rem:hadamard-state-independent-subtraction}
Let \(\omega\) and \(\omega'\) be Hadamard quasifree states for the same
operator \(P_M\) on the same globally hyperbolic spacetime.  Then
\[
  \omega_2-\omega'_2\in C^\infty(M\times M).
\]
\end{remark}

On each convex normal neighborhood, Definition~\ref{def:hadamard-quasifree-state}
gives
\[
  \omega_2=H_{\epsilon,\mu}+W_\omega,
  \qquad
  \omega'_2=H_{\epsilon,\mu}+W_{\omega'}
\]
after taking the distributional boundary value.  The same
\(H_{\epsilon,\mu}\) appears because \(U\) and \(V\) are fixed by
\((M,g)\), \(m\), and \(\xi\).  Therefore
\[
  \omega_2-\omega'_2=W_\omega-W_{\omega'}
\]
is smooth on that neighborhood.  Smoothness is local on \(M\times M\), so the
difference is smooth globally.

\section{Microlocal Form And Wick Products}

The local formula above is equivalent to the microlocal Hadamard condition
developed in Chapter~\ref{ch:curved-microlocal-spectrum-condition}:
\[
  \operatorname{WF}(\omega_2)
  =
  \{(x,k_x;y,-k_y): (x,k_x)\sim(y,k_y),\ k_x\in\overline V_+^*\}.
\]
Here \((x,k_x)\sim(y,k_y)\) means that \(x\) and \(y\) are connected by a
null geodesic and that \(k_y\) is the parallel transport of \(k_x\).  This
condition is the curved-spacetime replacement for positive energy support.

For a Hadamard state, define the Wick square by point splitting:
\[
  :\!\phi^2\!:_H(f)
  =
  \int_M f(x)
  \lim_{y\to x}
  \left(\phi(x)\phi(y)-H_{\epsilon,\mu}(x,y)\mathbf 1\right)
  \dd\operatorname{vol}_g(x).
\]
Remark~\ref{rem:hadamard-state-independent-subtraction} implies that
the subtraction removes the state-independent singularity.  The remaining
freedom is finite and local.  In four dimensions,
\[
  :\!\phi^2\!:'=
  :\!\phi^2\!:+c_1m^2+c_2R
\]
for constants \(c_1,c_2\), once covariance, scaling dimension, and the field
equation are imposed.  Stress-tensor point splitting in the next chapter is
the same construction with the bidifferential operator dictated by the
classical stress tensor.

\section{Renormalization Freedom Of The Stress Tensor}

The stress tensor is obtained by applying a bidifferential operator
\(D_{\mu\nu}(x,y)\) to the subtracted two-point expression and then taking
the diagonal limit:
\[
  \langle T_{\mu\nu}(x)\rangle_\omega
  =
  \lim_{y\to x}
  D_{\mu\nu}(x,y)
  \bigl(\omega_2(x,y)-H_{\epsilon,\mu}(x,y)\bigr)
  +
  C_{\mu\nu}[g](x).
\]
The tensor \(C_{\mu\nu}[g]\) is local and covariant in the background.  In
four dimensions, conservation and engineering dimension restrict it to a
linear combination of metric variations of
\[
  \int\sqrt{|g|},\qquad
  \int\sqrt{|g|}R,\qquad
  \int\sqrt{|g|}R^2,\qquad
  \int\sqrt{|g|}R_{\mu\nu}R^{\mu\nu},\qquad
  \int\sqrt{|g|}R_{\mu\nu\rho\sigma}R^{\mu\nu\rho\sigma},
\]
with constants specifying the finite local curvature coordinates of the
renormalized theory.  This statement is a covariance and conservation
classification of finite local terms; the explicit point-splitting
construction and the trace-anomaly consequences are developed in
Chapters~\ref{ch:point-splitting-stress-tensor-renormalization} and
\ref{ch:curved-trace-anomalies}.

\begin{openproblem}[Interacting curved-spacetime QFT]
\label{op:interacting-curved-spacetime-qft}
Construct interacting locally covariant QFTs with nonperturbative control of
state spaces, composite operators, stress tensors, and background dependence.
The construction should compare perturbative algebraic QFT, constructive
methods, lattice methods on discretized backgrounds, and Hamiltonian
approaches in examples where more than one construction exists.
\end{openproblem}
